\section{Short introduction}
More broadly speaking, control problems can be categorized
into regulation or tracking problems.
In regulation, the task is to stabilize the observed
system state to a set reference value by computing
error from the observed state to the desired reference,
and then use a control law to drive the error to zero.
In tracking problems, the task is to stabilize
the system state around a trajectory. Said another way,
there is a desired system state evolution,
which is simply a line in some high dimensional space, parametrized by time.
Thus, for each time instant, there is a desired system state.
The errors can be instantaneous, or integrated over a time horizon.
Usually when formulating a control law, the references' highest derivative 
is given as a feedforward term, while the others are used for 
feedback.

In many ways, robotics problems are just a subset of control
problems more broadly. This should be kept in mind,
especially since one can easily get distracted with all the geometry,
and the various competing control objectives.
Regulation problems are much less common compared
to tracking problems.
In this text, and in the code, regulation problems are referred to 
as ``\textbf{point to point motions}''.
Tracking problems also consist of generating the reference is some way.
Traditionally, this is referred to as
\textbf{motion planning}, and is
decomposed in the following steps:
\begin{itemize}
\item \textbf{path planning} - generating a sequence of waypoints
	(more common),
	or a smooth curve to a desired point
\item \textbf{interpolation} - fitting of a smooth curve
	to a set of planned waypoints
\item \textbf{trajectory generation} - construction of a
	time law on the smooth path, i.e. finding the mapping
	between the path parameter and time
\end{itemize}

In terms of manipulator control, the positional encoders
offer a massive benefit as they render the manipulator into
a fully observed system.
Otherwise, many challenges lie ahead.
They stem from the fact that there is no silver bullet control
scheme to solve every conceivable manipulator task at once.
There are many control objectives to satisfy:
\begin{itemize}
\item reaching the desired reference quickly and efficiently 
	(in terms of motor effort)
\item avoiding collisions of the robot with itself or with the environment
\item simultaneously tracking force and motion references
\item optimizing the robot's posture for the task at hand
\item ...
\end{itemize}
Even for a single manipulator, putting every objective
into one optimal control problem results in a nonlinear problem
which can not be solved for globaly.
Thus different control strategies, including motion planning,
need to employed depending on the situation. 

For example, one simple way to ensure safe interaction with the environment
is to use impedance control.
Namely, the manipulator may be equipped with force/torque sensors,
and the error dynamics can be turned into a spring-damper system.
This results in the robot being compliant, or ``soft'' in
response to an external force.

In case of mobile robots, things become substantially more complicated.
Localization is noisy, and often slow. The control strategy
needs to accommodate both of these features.
Additionally, in whole-body control tasks,
the task is often (very) underdetermined, leaving
many additional degrees of freedom that have to be accounted
for for good performance.


\subsection{Some common control approaches}
\subsubsection{Model-based}
If the task is to move the robot from one vector
of joint positions to another, or to follow a trajectory
of joint position, the task is fully defined
in \textbf{joint space}.
Control tasks for such problems are called
\textbf{joint space control} tasks.
However, we most often want to define tasks in terms of 
where we want the robot to be in the environment.
Purely geometric approaches to such problems
are referred to as 
\textbf{operational/Cartesian space control}.

Leveraging force/torque sensing can result
admitance control, impedance control or hybrid force-motion
control.

Optimal control problems in robots are usually formulated
in joint space, in the sense that the task is to solve for
a desired joint-space trajectory.

\subsubsection{Model-free (learning)}
All the hitherto mentioned control approaches are model-based,
meaning that they rely on the model of the robot
and the environment to compute the errors and solve for the
control input.
Recently, learning, and especially model-free learning,
gained a lot of popularity in robotics control.
It offers a huge benefit when there is no clear
way to turn robot motion requirements into control objectives,
and it very often simply beats model-based optimal control approaches.
This is because the learning algorithms can work on much 
sparser error signals, or in ML lingo, rewards.
Another huge benefit is avoiding detailed environment modelling,
which is a very hard problem in manipulation.
Not only are the models difficult to observe, but they often 
lead to NP-hard problems (planning for contacts in manipulation for example).
There are two big drawbacks however. The first is that 
learned controllers do not come with rigorously proven mathematical 
convergence and performance guarantees (and thus safety requirements)
The second is the so called ``sim-to-real'' gap.
Namely, learning takes many iterations, and is generally infeasible
on hardware for this reason. Thus the learning is done in simulation.
However, no efficient physics simulator corresponds to reality ---
approximations have to be made for the simulator to actually
be faster than a hardware experiment. The problem is that
the learning methods overfit to the simulation, and
are difficult to transfer into the real world.
This is a very active area of research, and there are interesting
solutions and avenues for work in terms of both problems.

\subsubsection{Which control approach to pick?}
Here it is important to consult the following evergreen 
rule of thumb in engineering: use the right tool for the job.
Since simplicity and explainability are always a good thing,
this means that one should use the simplest \textit{possible} control 
technique for the task at hand.
There is no point in learning inverse kinematics*,
just as much as there is no point in using inverse kinematics
and a FEM model of manipulate complex deformable objects.
